\chapter{Quantum inference}

\abstract{
    We here investigate, how inference on \ComputationActivationNetworks{} can be executed by quantum algorithms.
    While we find circuits corresponding to the computation networks, we exploit Grover operators for the activation of the computed qubits.
    We propose two inference algorithms:
    \begin{itemize}
        \item \textbf{Quantum rejection sampling:}
        When the canonical parameters of an \CompActNet{} are available, we prepare q-samples of the ancilla augmented \CompActNet{} and apply Grover rotations to achieve an square root advantage in sampling.
        \item \textbf{Quantum moment matching:}
        When the mean parameters to an \CompActNet{} are available, we iteratively construct the q-sample of an approximation.
    \end{itemize}
}

\section{Introduction}

\todo{
    We here present quantum inference schemes for CANs.
    All of them exploit computation circuits as a central building block and Grover operators as an activation mechanism.
    While for rejection sampling we add ancilla qubits, for moment estimation and matching we do not require further qubits.
}



% Classical TN contractions
Tensor network contractions have been identified as central inference mechanism of neuro-symbolic AI.
However, their exact execution is $\mathrm{NP}$-hard, and one either has to restrict itself to feasible situations, or apply approximations.


Tensor Network contractions can be performed by quantum computing as follows:
\begin{itemize}
    \item \textbf{Rejection Sampling based:}
    Can augment any tensor network by an Bayesian network, which is efficiently inferred by forward sampling.
    We prepare quantum states, which measurement statistics can be utilized to prepare samples from interesting distributions such as \ComputationActivationNetworks{}.
    The square root quantum advantage is obtained by amplitude amplification.
    \item \textbf{Phase Estimation based:}
    Measure the eigenvalues of a Grover operator to a formula and a q-sample.
    \item \textbf{Fourier Transform based:}
    Limited contractions (given promises on their possible values) can be formulated as Hidden Subgroup Problems and solved by Fourier transforms.
%    Here we investigate how we can exploit these as contraction provider for \tnreason{}.
    We are inspired by Deutsch-Jozsa algorithm, which we generalize here to a general moment computation problem.
\end{itemize}

%
Note that the rejection sampling scheme and the phase estimation are ignorant under the phase of the prepared state, and only the Fourier transform based scheme depend on the phase.

% NP Hardness of contraction:
Generic tensor network contraction is $\mathrm{NP}$-hard, reflecting the hardness of logical and probabilistic inference.
Quantum Computers can prepare states, which measurement distribution are specific network contractions, which can be hard to contract classically.
We study schemes to retrieve information about these states, for example by sampling schemes, and use the retrieved information in inference.

\subsection{Circuit Preparation schemes}

Circuit preparing schemes based on approximation:
\begin{itemize}
    \item Q-Alchemy \cite{araujo_low-rank_2023}, Q-Tucker \cite{blank_tucker_2026}
    \item Tensor-Network Optimization based (alternating schemes) \cite{rudolph_synergistic_2023,rudolph_decomposition_2023}
    \item Preparing a quantum circuit given an MPS requires logarithmic depth in the tensor order \cite{malz_preparation_2024}
\end{itemize}

Exact circuit preparing schemes for distributions:
\begin{itemize}
    \item Grover-Rudolph \cite{grover_creating_2002}, but no quantum advantage \cite{herbert_no_2021}
    \item Uniform controlled rotations \cite{mottonen_transformation_2005}
    \item Quantum Shannon decomposition \cite{shende_synthesis_2006}
\end{itemize}

Circuit simulation:
Since \qcreason{} can prepare quantum circuits to arbitrary tensor networks, it can also be used to simulate quantum circuits (with an overhead!).
\begin{itemize}
    \item \cite{sander_large-scale_2025,sander_quantum_2025}
\end{itemize}

\subsection{Related Works}

Quantum algorithms for linear algebra:
\begin{itemize}
    \item Contraction by SWAP and Hadamard test (see Appendix)
    \item Deutsch-Jozsa Algorithm
    \item HHL algorithm to solve linear equations \cite{harrow_quantum_2009}
\end{itemize}
%\section{Grover operator}\label{sec:groverOperator}

We derive quantum inference schemes based on the Grover operator, which we want to first introduce in full generality.

\subsection{Definition}

Let us assume we have two orthogonal quantum states $\qstateofat{\goodstatesymbol}{\shortcatvariables},\qstateofat{\badstatesymbol}{\shortcatvariables}$ spanning a subspace $\subspaceof{}=\spanof{\qstateofat{\goodstatesymbol}{\shortcatvariables},\qstateofat{\goodstatesymbol}{\shortcatvariables}}$ and let us choose another state
\begin{align*}
    \qstateofat{0}{\shortcatvariables}
    \in \subspaceof{}\, .
\end{align*}
Then there is $\rotanglesymbol\in[0,4\pi)$ with
\begin{align*}
    \qstateofat{0}{\shortcatvariables}
    = \sinof{\frac{\rotanglesymbol}{2}}\qstateofat{\goodstatesymbol}{\shortcatvariables}
    + \cosof{\frac{\rotanglesymbol}{2}}\qstateofat{\badstatesymbol}{\shortcatvariables} \, .
\end{align*}
Consider the reflection about the states $\qstateof{\badstatesymbol}$ and $\qstateof{0}$ by
\begin{align*}
    \reflectionofat{\badstatesymbol}{\inshortcatvariables,\outshortcatvariables}
    &= 2 \cdot \qstateofat{\badstatesymbol}{\inshortcatvariables} \otimes \qstateofat{\badstatesymbol}{\outshortcatvariables} - \identityat{\inshortcatvariables,\outshortcatvariables} \\
    \reflectionofat{0}{\inshortcatvariables,\outshortcatvariables}
    &= 2 \cdot \qstateofat{0}{\inshortcatvariables} \otimes \qstateofat{0}{\outshortcatvariables} - \identityat{\inshortcatvariables,\outshortcatvariables} \, ,
\end{align*}
and their concatenation
\begin{align*}
    \groverat{\inshortcatvariables,\outshortcatvariables}
    = \contractionof{
        \reflectionofat{\badstatesymbol}{\inshortcatvariables,\ccatvariableof{[\catorder]}{\midsymbol}},
        \reflectionofat{0}{\ccatvariableof{[\catorder]}{\midsymbol},\outshortcatvariables}
    }{\inshortcatvariables,\outshortcatvariables}
\end{align*}
which is called the \emph{Grover operator}.

\subsection{Action on projected quantum states}

The Grover operator leaves any quantum state $\qstate$ invariant, which is orthogonal to $\qstateof{\goodstatesymbol}$ and $\qstateof{\badstatesymbol}$ (such states get a factor $(-1)$ by both reflection operations).
The projection of the quantum state onto $\subspaceof{}$ gets rotated by the Grover operator (see \figref{fig:rotationSketch}).
To be more precise, we denote the projection by $(\mathcal{P}\qstate)[\selvariable]$ which has the coordinates
\begin{align*}
    \left(\mathcal{P}_{\subspaceof{}}\qstate\right)\left[\selvariable=0\right]
    &= \contraction{\left(\qstatewith\right)^{\dagger},\qstateofat{\goodstatesymbol}{\shortcatvariables}}\\
    \left(\mathcal{P}_{\subspaceof{}}\qstate\right)\left[\selvariable=1\right]
    &= \contraction{\left(\qstatewith\right)^{\dagger},\qstateofat{\goodstatesymbol}{\shortcatvariables}} \, .
\end{align*}
We define a matrix representation of the Grover operator as
\begin{align*}
    \groverofat{\mathrm{red}}{\cselvariable{\outsymbol},\cselvariable{\insymbol}}
    = \coloredmatrixof{
        \cosof{\rotanglesymbol} & -\sinof{\rotanglesymbol} \\
        \sinof{\rotanglesymbol} & \cosof{\rotanglesymbol}
    }{\cselvariable{\outsymbol},\cselvariable{\insymbol}}
\end{align*}
and it holds for any quantum state $\qstatewith$ that
\begin{align*}
    \left(\mathcal{P}_{\subspaceof{}}\groversymbol\qstate\right)\left[\cselvariable{\outsymbol}\right]
    = \contractionof{\groverofat{\mathrm{red}}{\cselvariable{\outsymbol},\cselvariable{\insymbol}},
        \left(\mathcal{P}_{\subspaceof{}}\qstate\right)\left[\cselvariable{\insymbol}\right]
    }{\outsymbol} \, .
\end{align*}
The Grover operator therefore has the eigenvalues:
\begin{itemize}
    \item $\expof{\pm i \rotanglesymbol}$ with the eigen spaces spanned by $\qstateof{\goodstatesymbol}\pm i \qstateof{\badstatesymbol}$
    \item $1$ with the eigen space by the orthogonal complement of $\subspaceof{}$
\end{itemize}
Application the Grover operator $\repnum\in\nn$ times on $\qstateof{0}$ results in the state
\begin{align*}
    \qstateofat{\repnum}{\shortcatvariables}
    = \sinof{\left(\frac{1}{2}+\repnum\right)\rotanglesymbol} \qstateofat{\goodstatesymbol}{\shortcatvariables}
    + \cosof{\left(\frac{1}{2}+\repnum\right)\rotanglesymbol} \qstateofat{\badstatesymbol}{\shortcatvariables} \, .
\end{align*}

\begin{figure}[t]
    \begin{center}
        \begin{tikzpicture}[>=stealth, scale=4]

    % Coordinates
    \coordinate (O) at (0,0);
    \coordinate (W) at (0,1);     % Target state |w>
    \coordinate (Sprime) at (1,0); % All other states |s'>

    % Draw Axes
    \draw[->, thick] (-0.1,0) -- (1.1,0) node[right] {\corelabelsize $\qstateof{\badstatesymbol}$};
    \draw[->, thick] (0,-0.1) -- (0,1.1) node[above] {\corelabelsize $\qstateof{\goodstatesymbol}$};

    \def\sthalfangle{15}
    \def\nextangle{3*\sthalfangle} % 3 * theta
    \draw[->, thick] (O) -- (\sthalfangle:0.9) node[right] {\corelabelsize $\qstateof{0}=\sinof{\frac{\rotanglesymbol}{2}}\qstateof{\goodstatesymbol}+\cosof{\frac{\rotanglesymbol}{2}}\qstateof{\badstatesymbol}$};
    \draw[dashed, thick] (0.4,0) arc (0:\sthalfangle:0.4) node[midway, right, xshift=2pt] {\corelabelsize $\frac{\rotanglesymbol}{2}$};

    \draw[<-, thick, blue] (-1*\sthalfangle:0.6) arc (-\sthalfangle:\sthalfangle:0.6) node[midway, right, yshift=-5pt, xshift=2pt] {\corelabelsize $\rotanglesymbol$};

    % Draw the Oracle Reflection (reflected across s' axis)
    \draw[->, thick, blue] (O) -- (-\sthalfangle:0.9) node[right] {\corelabelsize $-\sinof{\frac{\rotanglesymbol}{2}}\qstateof{\goodstatesymbol}+\cosof{\frac{\rotanglesymbol}{2}}\qstateof{\badstatesymbol}$};


    \draw[->, thick, red] (-1*\sthalfangle:0.8) arc (-\sthalfangle:3*\sthalfangle:0.8) node[midway, right, yshift=15pt, xshift=-5pt] {\corelabelsize $2\rotanglesymbol$};

    % Draw the first Grover Iteration (rotated by 2*theta from original |s>)

    \draw[->, red, thick] (O) -- (\nextangle:0.9) node[above right] {\corelabelsize $\qstateof{1}=\sinof{\left(\frac{1}{2}+1\right)\rotanglesymbol}\qstateof{\goodstatesymbol}+\cosof{\left(\frac{1}{2}+1\right)\rotanglesymbol}\qstateof{\badstatesymbol}$};

    \node[below left] at (O) {\corelabelsize $0$};

\end{tikzpicture}
    \end{center}
    \caption{Grover operator as a rotation in $\spanof{\qstateof{\goodstatesymbol},\qstateof{\badstatesymbol}}$.}
    \label{fig:rotationSketch}
\end{figure}

\subsection{Circuit Realization of the reflection operators}

Let us now investigate how the reflection operations can be realized by circuits.


% THIS WOULD MAKE A DIFFERENCE IN QUANTUM PHASE ESTIMATION!

\subsubsection{Construction based on controlled Pauli-Z}

First, the reflection about the state $\onehotmapofat{\onetuple{[\catorder]}}{\shortcatvariables}$ is realized by a multiple controlled Pauli-Z (see \figref{fig:reflectionZRealization})
%    \begin{align*}
%        -\reflectionofat{\goodstatesymbol}{\inshortcatvariables,\outshortcatvariables}
%        = \mathrm{MC}^{\seldim-1}\paulizat{\cavariableof{\seldim-1}{\insymbol},\cavariableof{\seldim-1}{\outsymbol},\avariableof{[\seldim-1]}}
%    \end{align*}
\begin{align*}
    -\reflectionofat{\onehotmapof{\onetuple{[\catorder]}}}{\inshortcatvariables,\outshortcatvariables}
    = \contractionof{
        \mathrm{MC^{\catorder-1}}\paulizat{\ccatvariableof{\catorder-1}{\insymbol},\ccatvariableof{\catorder-1}{\outsymbol},\catvariableof{[\catorder-1]}},
        \identityat{\ccatvariableof{[\catorder-1]},\ccatvariableof{[\catorder-1]}{\insymbol},\ccatvariableof{[\catorder-1]}{\outsymbol}}
    }{\inshortcatvariables,\outshortcatvariables}
\end{align*}
Note that a global phase of $-1$ is introduced, which leaves the measurement distribution invariant.
% Extending to e_1
Adding layers of Pauli-X gates to each $\catenumeratorin$ before and after $-\reflectionofat{\onehotmapof{\onetuple{[\catorder]}}}{\inshortcatvariables,\outshortcatvariables}$ (see \figref{fig:reflectionZRealization}) prepares the negated reflection $-\reflectionofat{\onehotmapof{\zerotuple{[\catorder]}}}{\inshortcatvariables,\outshortcatvariables}$, since
\begin{align*}
    \tbasisat{\ccatvariableof{\catenumerator}{\outsymbol}}
    = \contractionof{
        \paulizat{\ccatvariableof{\catenumerator}{\outsymbol},\ccatvariableof{\catenumerator}{\insymbol}},\fbasisat{\ccatvariableof{\catenumerator}{\insymbol}}
    }{\ccatvariableof{\catenumerator}{\outsymbol}} \, .
\end{align*}
% Extending to U
Having a unitary $\unitarysymbol$ preparing $\qstate$ by action on the ground state, we realize the reflection (see \figref{fig:reflectionZRealization})
\begin{align*}
    \reflectionofat{\qstate}{\inshortcatvariables,\outshortcatvariables}
    = \contractionof{
        \unitaryofat{\dagger}{\inshortcatvariables,\ccatvariableof{[\catorder]}{\midsymbol,0}},
        \reflectionofat{\onehotmapof{\zerotuple{[\catorder]}}}{\ccatvariableof{[\catorder]}{\midsymbol,0},\ccatvariableof{[\catorder]}{\midsymbol,1}},
        \unitaryat{\ccatvariableof{[\catorder]}{\midsymbol,1},\outshortcatvariables}
    }{\inshortcatvariables,\outshortcatvariables} \, .
\end{align*}


\begin{figure}[t]
    \begin{center}
        \begin{tikzpicture}[scale=0.35, thick] % , baseline = -3.5pt

%    \drawsmalltensor{0}{9}{\corelabelsize $\fbasis$}
%    \draw (1,9) -- (2,9);
%    \drawsmalltensor{3}{9}{\corelabelsize $\paulixsymbol$}
%    \draw (4,9) -- (5,9);
%    \drawsmalltensor{6}{9}{\corelabelsize $\hgate$}
%
%    %\draw (-1,5) rectangle (1,7);
%    \drawsmalltensor{0}{6}{\corelabelsize $\fbasis$}
%    \draw (1,6) -- (2,6);
%    \drawsmalltensor{0}{2}{\corelabelsize $\fbasis$}
%    \draw (1,2) -- (2,2);
%    \draw (2,1) rectangle (4,7);
%    \drawtextnode{3}{4}{\corelabelsize $\unitarysymbol$}
%    \draw (4,6) -- (12,6);
%    \draw (4,2) -- (12,2);
%
%    %\draw (2,1) rectangle (4,3);
%    %\drawtextnode{3}{2}{\corelabelsize $\hgate$}
%
%
%    \drawtextnode{-3}{9}{\corelabelsize $\headvariable$}
%    \draw[dashed] (-2,7.5) -- (30.5,7.5);
%    \drawtextnode{-3}{4}{\corelabelsize $\shortcatvariables$}
%    \drawtextnode{0}{4.5}{\corelabelsize $\vdots$}
%
%    \draw[dashed] (7.5,12) -- (7.5,-1) node[below] {\corelabelsize $\qstateofat{0}{\shortcatvariables} \otimes \hminusbasat{\headvariable}$};
%
%    \begin{scope}[shift={(4,0)}]
%        \draw (3,9) -- (4,9);
%        \draw (4,8) rectangle (7,10);
%        \node[anchor=center] (text) at (5.5,9) {\corelabelsize $\qcbencodingof{\exformula}$};
%        \draw (5,6) -- (5,8);
%        \drawvariabledot{5}{6}
%        \draw (6,2) -- (6,8);
%        \drawvariabledot{6}{2}
%    \end{scope}
%
%    \drawtextnode{9.5}{11}{\corelabelsize $\reflectionof{\badstatesymbol}$}
%
%    \draw[dashed] (11.5,11) -- (11.5,0.5);

    \drawtextnode{9.5}{4}{\corelabelsize $\inshortcatvariables$}
    \draw (10,6) -- (12,6);
    \drawtextnode{11.5}{4.25}{$\vdots$}
    \draw (10,2) -- (12,2);
    \draw (12,1) rectangle (14,7);
    \drawtextnode{13}{4}{\corelabelsize $\unitaryof{\dagger}$}
    \draw (14,6) -- (15,6);
    \drawsmalltensor{16}{6}{\corelabelsize $\paulixsymbol$}
    \draw (14,2) -- (15,2);
    \drawsmalltensor{16}{2}{\corelabelsize $\paulixsymbol$}
    \draw (17,2) -- (21,2);
    \draw (17,6) -- (18,6);
    \drawsmalltensor{19}{6}{\corelabelsize $\paulizsymbol$}
    \draw (18.5,5) -- (18.5,4);
    \drawvariabledot{18.5}{4}
    \node[rotate=-60] at (18.9,2.75) {$\cdots$};
    \draw (19.5,5) -- (19.5,2);
    \drawvariabledot{19.5}{2}

    \draw (20,6) -- (21,6);
    \drawsmalltensor{22}{6}{\corelabelsize $\paulixsymbol$}
    \draw (23,6) -- (24,6);
    \drawsmalltensor{22}{2}{\corelabelsize $\paulixsymbol$}
    \draw (23,2) -- (24,2);
    \draw (24,1) rectangle (26,7);
    \drawtextnode{25}{4}{\corelabelsize $\unitarysymbol$}
    \draw (26,6) -- (28,6);
    \draw (26,2) -- (28,2);
    \drawtextnode{27}{4.25}{$\vdots$}
    \drawtextnode{28.5}{4}{\corelabelsize $\outshortcatvariables$}
    %\draw (11,9)--(28,9);

%    \drawtextnode{19}{11}{\corelabelsize $\reflectionof{0}$}

    \draw[dashed] (23.5,0) -- (23.5,9.25) -- (14.5,9.25) -- (14.5,0) -- cycle;
    \drawtextnode{16}{8.5}{\corelabelsize $-\reflectionof{\onehotmapof{\zerotuple{[\catorder]}}}$}

    \draw[dashed] (20.5,1.25) -- (20.5,8.75) -- (17.5,8.75) -- (17.5,1.25) -- cycle;
    \drawtextnode{19}{8}{\corelabelsize $-\reflectionof{\onehotmapof{\onetuple{[\catorder]}}}$}

    \begin{scope}[shift={(-12,0)}]
        \drawtextnode{9.5}{4}{\corelabelsize $\inshortcatvariables$}
        \draw (10,6) -- (12,6);
        \drawtextnode{11.5}{4.25}{$\vdots$}
        \draw (10,2) -- (12,2);
        \draw (12,1) rectangle (14,7);
        \drawtextnode{13}{4}{\corelabelsize $-\reflectionof{\qstate}$}
        \draw (14,6) -- (15,6);

        \draw (14,6) -- (16,6);
        \drawtextnode{14.5}{4.25}{$\vdots$}
        \draw (14,2) -- (16,2);
        \drawtextnode{16.5}{4}{\corelabelsize $\outshortcatvariables$}

        \drawtextnode{19}{4}{${=}$}
    \end{scope}


    %\draw[dashed]  -- ;
    %\draw[dashed]  -- ; % node[below] {\corelabelsize $\qstateofat{1}{\shortcatvariables} \otimes \hminusbasat{\headvariable}$};


\end{tikzpicture}
    \end{center}
    \caption{
        Preparation of the negated reflection operator $-\reflectionof{\qstate}$ about a state $\qstate$, which is prepared by a unitary $\unitarysymbol$ acting on the ground state $\onehotmapofat{\zerotuple{[\catorder]}}{\shortcatvariables}$.
        The multiply controlled Pauli-Z is the negated reflection $-\reflectionof{\onetuple{[\catorder]}}$, and extended to the negated reflection $-\reflectionof{\onetuple{[\catorder]}}$ about the ground state by adding layers of single qubit Pauli-X.
        When concatenating with $\unitaryof{\dagger}$ and $\unitarysymbol$, the $-\reflectionof{\qstate}$ results.
    }
    \label{fig:reflectionZRealization}
\end{figure}

\subsubsection{Construction based on oracle with ancilla head qubit}

Alternatively we can use the phasekick technique when acting on an disentangled ancilla qubit prepared in the Hadamard basis state
\begin{align*}
    \hminusbasat{\headvariable}
    = \frac{1}{\sqrt{2}} \cdot \coloredmatrixof{1 \\ -1}{\headvariable} \, .
\end{align*}
We notice, that
\begin{align*}
    \contractionof{
        \qcbencodingofat{\exformula}{\cheadvariable{\outsymbol},\cheadvariable{\insymbol},\indexedshortcatvariables},
        \hminusbasat{\cheadvariable{\insymbol}}
    }{\cheadvariable{\outsymbol}}
    = \begin{cases}
          \hminusbasat{\cheadvariable{\insymbol}} & \ifspace \exformula(\shortcatindices) = 0\\
          -\hminusbasat{\cheadvariable{\insymbol}} & \ifspace \exformula(\shortcatindices) = 1
    \end{cases}
\end{align*}

Let us consider the span of one-hot encoded models of $\lnot\exformula$, namely
\begin{align*}
    \subspaceof{\lnot\exformula}
    \coloneqq \spanof{\onehotmapofat{\shortcatindices}{\shortcatvariables}\wcols\exformula(\shortcatindices)=0}
\end{align*}
and the span of one-hot encoded models
\begin{align*}
    \subspaceof{\exformula}
    \coloneqq\spanof{\onehotmapofat{\shortcatindices}{\shortcatvariables}\wcols\exformula(\shortcatindices)=1} \, .
\end{align*}
The reflection across $\subspaceof{\lnot\exformula}$ is then performed by (see \figref{fig:reflectionCCRealization})
\begin{align*}
    \reflectionofat{\subspaceof{\lnot\exformula}}{\outshortcatvariables,\inshortcatvariables}
    = \contractionof{
        \hminusbasat{\cheadvariable{\outsymbol}},
        \qcbencodingofat{\exformula}{\cheadvariable{\outsymbol},\cheadvariable{\insymbol},\shortcatvariables},
        \hminusbasat{\cheadvariable{\insymbol}},
        \identityat{\shortcatvariables,\outshortcatvariables,\inshortcatvariables}
    }{\outshortcatvariables,\inshortcatvariables}
\end{align*}
We further have
\begin{align*}
    \reflectionofat{\subspaceof{\exformula}}{\outshortcatvariables,\inshortcatvariables}
    = -\reflectionofat{\left(\subspaceof{\exformula}\right)^{\perp}}{\outshortcatvariables,\inshortcatvariables}
    = -\reflectionofat{\subspaceof{\lnot\exformula}}{\outshortcatvariables,\inshortcatvariables} \, .
\end{align*}
When preparing a disentangled ancilla qubit $\headvariable$ in the basis state $\hminusbasat{\headvariable}$, the oracle $\qcbencodingof{\exformula}$ is effectively a negated reflection across the span of one-hot encoded models of $\exformula$.

%This is an effective
%This further coincides with the effective negative reflection across the span of one-hot encoded models
%\begin{align*}
%    \subspaceof{\exformula}
%    \coloneqq\spanof{\onehotmapofat{\shortcatindices}{\shortcatvariables}\wcols\exformula(\shortcatindices)=1} \, .
%\end{align*}

%We will investigate different effective schemes to realize $\reflectionof{\badstatesymbol}$ later.
%\begin{itemize}
%    \item When $\qstateof{\goodstatesymbol}=\onehotmapofat{\onetuple{[\seldim]}}{\avariables}\otimes $ and $\qstate$
%    \item Reflection across the hyperplane $\spanof{\onehotmapofat{\shortcatindices}{\shortcatvariables}\wcols\exformula(\shortcatindices)=0}$ effectively by $\qcbencodingof{\exformula}$ acting on an auxiliary disentangled qubit in $\hminusbasat{\headvariable}$
%\end{itemize}

\begin{figure}[t]
    \begin{center}
        \begin{tikzpicture}[scale=0.35, thick] % , baseline = -3.5pt


    %\drawtextnode{9.5}{4}{\corelabelsize $\inshortcatvariables$}

    \drawsmalltensor{10}{8}{\corelabelsize $\hminusbas$}
    \draw (11,8) -- (12,8);
    \draw (12,7) rectangle (15,9);
    \drawtextnode{13.5}{8}{\corelabelsize $\qcbencodingof{\formula}$}
    \draw (13,7) -- (13,6);
    \drawvariabledot{13}{6}
    \draw (14,7) -- (14,2);
    \drawvariabledot{14}{2}
    \draw (15,8) -- (16,8);
    \drawsmalltensor{17}{8}{\corelabelsize $\hminusbas$}

    \drawtextnode{10.5}{4}{\corelabelsize $\inshortcatvariables$}
    \draw (11,6) -- (16,6);
    \drawtextnode{12.5}{4.25}{$\vdots$}
    \draw (11,2) -- (16,2);
    \drawtextnode{16.5}{4}{\corelabelsize $\outshortcatvariables$}
    %\draw (12,1) rectangle (14,7);
    %\drawtextnode{13}{4}{\corelabelsize $\unitaryof{\dagger}$}
    %\draw (14,6) -- (15,6);

    \begin{scope}[shift={(-12,0)}]
        \drawtextnode{9.5}{4}{\corelabelsize $\inshortcatvariables$}
        \draw (10,6) -- (12,6);
        \drawtextnode{11.5}{4.25}{$\vdots$}
        \draw (10,2) -- (12,2);
        \draw (12,1) rectangle (14,7);
        \drawtextnode{13}{4}{\corelabelsize $-\reflectionof{\subspaceof{\formula}}$}
        \draw (14,6) -- (15,6);

        \draw (14,6) -- (16,6);
        \drawtextnode{14.5}{4.25}{$\vdots$}
        \draw (14,2) -- (16,2);
        \drawtextnode{16.5}{4}{\corelabelsize $\outshortcatvariables$}

        \drawtextnode{19}{4}{${=}$}
    \end{scope}


    %\draw[dashed]  -- ;
    %\draw[dashed]  -- ; % node[below] {\corelabelsize $\qstateofat{1}{\shortcatvariables} \otimes \hminusbasat{\headvariable}$};


\end{tikzpicture}
    \end{center}
    \caption{
        Realization of the reflection across the subspace of one-hot encoded models to a Boolean function $\formula$.
        An ancilla qubit is prepared in the Hadamard basis state $\hminusbas$ and serves as the head qubit of the oracle $\qcbencodingof{\formula}$.
        By the phase kickback effect, the qubit stays disentangled and the reflection $-\reflectionof{\subspaceof{\formula}}$ is effectively performed on the qubits $\shortcatvariables$.
    }
    \label{fig:reflectionCCRealization}
\end{figure}



\subimport{rejection-sampling/}{rejection-sampling}
\subimport{moment-amplification/}{moment-amplification}

%\section{Quantum Rejection Sampling}\label{sec:sampling}

We present the Quantum Rejection Sampling method in two steps:
\begin{itemize}
    \item {\bf Ancilla augmentation:} Find an ancilla augmentation of a distribution, such that there is a unitary $\unitarysymbol$ preparing a state with the augmented distribution as measurement distribution
    \item {\bf Acceptance amplification:} Amplify the amplitude of the ancilla qubits in state $1$ by repeated application of the Grover operator
\end{itemize}
Compared with the rejection sampling approach on the start augmentation of the distribution, we achieve a square-root speed-up in the expected number of gates to be applied per sample.

We in particular apply this scheme to sample from \ComputationActivationNetworks{}, where we apply the q-sample preparation scheme of \secref{sec:canQSample}.

\subsection{Acceptance Amplification}

Let there be a quantum state, which measurement distribution is an ancilla augmentation $\secprobat{\avariables,\shortcatvariables}$ of a target distribution $\probwith$.
We now show how Grover operators can be used to amplify the acceptance state and improve the efficiency of the rejection sampling method through an increase of the acceptance probability.
This is an instance of amplitude amplification.

\subsubsection{Grover construction}

We relate the state $\qstateof{\goodstatesymbol}$ with the event $\avariables=\onetuple{[\seldim]}$, by constructing it as
\begin{align*}
    \qstateofat{\goodstatesymbol}{\avariables,\headvariables,\shortcatvariables}
    =\onehotmapofat{\onetuple{[\seldim]}}{\avariables}
    \otimes \contractionof{\bencodingofat{\sstat}{\headvariables,\shortcatvariables},\sqrt{\probwith}}{\headvariables,\shortcatvariables}
\end{align*}
The complement state $\qstateof{\badstatesymbol}$ is constructed if $\acceptanceprob\neq1$ (in which case the acceptance amplilification scheme would not be necessary)
\begin{align*}
    \probofat{\badstatesymbol}{\avariables,\headvariables,\shortcatvariables}
    = \frac{1}{1-\acceptanceprob}
    \cdot\left(
             \secprobat{\avariables,\headvariables,\shortcatvariables}-\acceptanceprob \cdot \onehotmapofat{\onetuple{[\seldim]}}{\avariables}\otimes\probat{\headvariables,\shortcatvariables}
    \right)
\end{align*}
and
\begin{align*}
    \qstateofat{\badstatesymbol}{\avariables,\headvariables,\shortcatvariables}=\sqrt{\probofat{\perp}{\avariables,\headvariables,\shortcatvariables}} \, .
\end{align*}

% Grover operator instance
%We build the Grover operator (see \secref{sec:groverOperator}) with respect to
%\begin{itemize}
%%    \item $$
%    \item $\qstateof{\badstatesymbol}=\sqrt{\probof{\perp}}$
%\end{itemize}
%and the state

We thus have
\begin{align*}
    \qstateofat{0}{\avariables,\headvariables,\shortcatvariables}
    = \sinof{\frac{\rotanglesymbol}{2}} \qstateofat{0}{\avariables,\headvariables,\shortcatvariables}
    + \cosof{\frac{\rotanglesymbol}{2}} \qstateofat{1}{\avariables,\headvariables,\shortcatvariables}
\end{align*}
where the $\rotanglesymbol$ is related to the acceptance probability as
\begin{align*}
    \sinof{\frac{\rotanglesymbol}{2}}
    = \sqrt{\acceptanceprob} \, .
\end{align*}



\subsubsection{Circuit Realization}

Both reflection operators are realized as their negations using the Pauli-Z scheme (see \secref{sec:groverOperator}).
The negations chancel in their concatenation, which is therefore the Grover operator.


\subsubsection{Ancilla augmentation preserved}

The measurement distribution to any quantum state in $\subspaceof{}\backslash\spanof{\qstateof{\badstatesymbol}}$ is an ancilla augmentation of $\probwith$. % Need to show?
When amplifying the acceptance probability we stay in the subspace $\subspaceof{}$.
When we do not overrotate and since we start with $\acceptanceprob>0$ the amplified state is never in $\spanof{\qstateof{\badstatesymbol}}$.

%% Realization of the grover operator
%While the reflection $\reflectionof{0}$ is realized as described in \secref{sec:groverOperator}, the reflection $\reflectionof{\badstatesymbol}$ is a multi-controlled Pauli-Z on the ancilla variables. \
%\red{No, we use the negated about the good state!}



\subsection{Optimal rotation numbers}

% Application
Now the application of $\repnum$ grover operators $\groversymbol$ on $\qstateof{0}$ prepares a state
\begin{align*}
    \qstateofat{\repnum}{\avariables,\shortcatvariables}
    = \sinof{\left(\frac{1}{2}+\repnum\right)\rotanglesymbol} \qstateofat{\goodstatesymbol}{\avariables,\shortcatvariables}
    + \cosof{\left(\frac{1}{2}+\repnum\right)\rotanglesymbol} \qstateofat{\badstatesymbol}{\avariables,\shortcatvariables}
\end{align*}
The state is closest to $\qstateof{\goodstatesymbol}$ when
\begin{align*}
    \left(\frac{1}{2}+\repnum\right)
    \rotanglesymbol \approx \frac{\pi}{2} \, .
\end{align*}

Estimating for small $\acceptanceprob$ (tight for small $\acceptanceprob$)
\begin{align*}
    \sqrt{\acceptanceprob} = \sinof{\frac{\rotanglesymbol}{2}}\leq \frac{\rotanglesymbol}{2}
\end{align*}
we get a bound
\begin{align*}
    \repnumof{\mathrm{opt}} \leq \frac{\pi}{4\sqrt{\acceptanceprob}} \, .
\end{align*}
This shows that we have a square root improvement on classical rejection sampling, which would require $\frac{1}{\acceptanceprob}$ repetitions to produce a sample.
Note however that we compare quantum circuit length with the number of classical repetitions in sampling.


%Note, that the variable qubits are uniformly distributed when only the computation circuit is applied.
%When sampling the probability distribution, we need the ancilla qubits to be in state $1$ in order for the sample to be valid.
%Any other states will have to be rejected.

%\textbf{Open Question:} Is there a way to avoid amplitude amplification and use a more direct circuit implementation of the activation network?
%- Cannot be the case, when the encoding is determined by the activation tensor alone: Needs to use the computated statistic as well.
%Negative result in \cite[Section~6.6]{nielsen_quantum_2010}: When using multiple applications of \computationCircuit{} in combination with further unitaries, need at least as many applications as with amplitude amplification.

\subsection{Applications on \CompActNets{}}

The above scheme prepares the ancilla augmented \ComputationActivationNetwork{}
\begin{align*}
    \secprobat{\avariables,\headvariables,\shortcatvariables}
    = \acceptanceprob \cdot \onehotmapofat{\onetuple{[\seldim]}}{\avariables}
    \otimes \contractionof{\bencodingofat{\sstat}{\headvariables,\shortcatvariables},\probwith}{\headvariables,\shortcatvariables}
    + (1-\acceptanceprob) \probofat{\perp}{\shortcatvariables,\headvariables,\avariables}
\end{align*}
and the prepared state is the combination of q-samples
\begin{align*}
    \qstateofat{0}{\avariables,\headvariables,\shortcatvariables}
    = \sqrt{\acceptanceprob} \cdot \onehotmapofat{\onetuple{[\seldim]}}{\avariables}
    \otimes \contractionof{\bencodingofat{\sstat}{\headvariables,\shortcatvariables},\sqrt{\probwith}}{\headvariables,\shortcatvariables}
    + \sqrt{(1-\acceptanceprob)} \sqrt{\probofat{\perp}{\shortcatvariables,\headvariables,\avariables}} \, .
\end{align*}
%Here we understand $\headvariables$ as extending $\shortcatvariables$.

Sampling from the augmented \ComputationActivationNetwork{} can be also done classically by forward sampling:
Just draw the variables from uniform, calculate the value qubit by a logical circuit inference and accept with probability by the computed value.
To achieve a quantum advantage we manipulate the circuit first by the amplitude amplification scheme and increase the acceptance probability.

\subsection{Examples}

%% To be included:

%\begin{example}[Sample a model of a propositional formula]

    Acceptance rate:
    \begin{align*}
        \acceptanceprob
        = \frac{\contraction{\formula}}{2^{\catorder}}
    \end{align*}

    Angle:
    \begin{align*}
        \rotanglesymbol
        = 2 \cdot arcsin\left(\sqrt{\frac{\contraction{\formula}}{2^{\catorder}}}\right)
    \end{align*}

\end{example}
%\begin{example}[Sampling from the toy accounting]

    For positive $\canparam$: $\repnumof{\mathrm{opt}}=0$

    For large negative $\canparam$: Effectively only one model (i.e. minimal acceptance prob), $\repnumof{\mathrm{opt}}=1$

\end{example}



%\section{Moment Amplification}

% Comparison with Moment Estimation
Instead of measuring the phases of the effective Grover operator, we now investigate the effect of uncontrolled Grover operations on the distributions.
We find an analogon in discrete increases of the canonical parameters when the measurement distribution of the start state is in an exponential family.


%We here apply \computationCircuits{} on anti-symmetric ancilla states (as for sign encoding)
%\begin{itemize}
%    \item A single \computationCircuit{} effectively prepares a reflection across the models of the computed formula
%    \item Building the reflections based on ancilla ground states (as for basis encoding), we would need an additional Pauli-Z and two \computationCircuits{}
%    \item Further ancilla qubits used in circuit decompositions need to be uncomputed (i.e. by applying the same partial \computationCircuit{} again) for further usage
%\end{itemize}

\subsection{Interpretation of Grover rotations}

%\subsection{Rotation along a formula}

Rotating $\repnum$ times across the initial state then gives the tensor product of the antisymmetric ancilla state with
\begin{align*}
    \qstateofat{\repnum}{\shortcatvariables}
    = \sinof{\left(\frac{1}{2}+\repnum\right)\rotanglesymbol} \qstateofat{\probof{0},\exformula}{\shortcatvariables} +
    \cosof{\left(\frac{1}{2}+\repnum\right)\rotanglesymbol} \qstateofat{\probof{0},\lnot\exformula}{\shortcatvariables} \, .
\end{align*}
This reflection can be performed by the operator on the distributed qubits
\begin{align*}
    U^t \left(
            \identityat{\shortcatvariables^{\insymbol},\shortcatvariables^{\outsymbol}}-2\onehotmapofat{0}{\shortcatvariables^{\insymbol}} \otimes \onehotmapofat{0}{\shortcatvariables^{\outsymbol}}
    \right) U \, .
\end{align*}


\subsubsection{Mean parameter interpretation}
% Mean parameter interpretation

% As exponential family parameters
The mean parameter of $\exformula$ with respect to the measurement distribution is
\begin{align*}
    \meanparam^{\repnum}
    = \contraction{\probofat{\repnum}{\shortcatvariables},\formulawith}
    = \sinof{\left(\frac{1}{2}+\repnum\right)\rotanglesymbol}^2
    = \sinof{\left(1+2\repnum\right)\sin^{-1}\left(\sqrt{\meanparam^{0}}\right)}^2
    %= \left(\frac{
    %    2\cdot \sinof{\left(\frac{1}{2}+\repnum\right)\rotanglesymbol} \cdot \sin^{-1}\left(\sqrt{\meanparam^{0}}\right)
    %}{\meanparam^{0}}\right)^2
\end{align*}
Here we used that $\rotanglesymbol = 2\cdot\sin^{-1}\left(\sqrt{\meanparam^{0}}\right)$.


\subsubsection{Canonical parameter interpretation}

%\begin{figure}[t]
%    \begin{center}
%        \input{tikz/moment-amplification/rotation-sketch}
%    \end{center}
%    \caption{Geometric visualization of the states $\qstateof{1}$ resulting from a Grover rotation on $\qstateof{0}$, in the subspace spanned by $\qstateof{\probof{0},\exformula}$ and $\qstateof{\probof{0},\lnot\exformula}$.
%    The effective reflection by the \computationCircuit{} results in the gray state $\qstateof{\badstatesymbol}-\qstateof{\goodstatesymbol}$.
%    Reflecting again on the initial state $\qstateof{0}$ results in the amplified state $\qstateof{1}$.
%    Each such Grover iteration is therefore a rotation by the angle $\rotanglesymbol$ in the.}
%\end{figure}



\begin{lemma}
    \label{lem:rotationCanParamIncrease}
    \red{Works also for more generic measurement distributions!}
    Let $\qstateof{0}$ be the q-sample of the elementary \ComputationActivationNetwork{} with canonical parameters $\hybridparam$, and let us amplify the formula $\formulaof{\selindex}$.
    Then $\qstateof{\repnum}$ is the q-sample of the elementary \ComputationActivationNetwork{} with canonical parameters
    \begin{align*}
        \hardlegset^{\repnum} =
        \begin{cases}
            \hardlegset \cup \{\selindex\}, \headindexof{\hardlegset}^{\selindex}
            = (\headindexof{\hardlegset},1) & \ifspace \sinof{(1+2\cdot\repnum)\sin^{-1}\left(\sqrt{\meanparamofat{0}{\indexedselvariable}}\right)} = 1 \\
            \hardlegset, \headindexof{\hardlegset} & \elsetext
        \end{cases}
    \end{align*}
    and
    \begin{align*}
        \canparamofat{\repnum}{\indexedselvariable} =
        \begin{cases}
            0 & \ifspace  \sinof{(1+2\cdot\repnum)\sin^{-1}\left(\sqrt{\meanparamofat{0}{\indexedselvariable}}\right)} = 1 \\
            \canparamat{\indexedselvariable} + \frac{
                \cos^2\left(\frac{\rotanglesymbol}{2}\right) \left(1-
                                                                 \cos^2\left(\left(\frac{1}{2}+\repnum\right) \rotanglesymbol \right)
                \right)
            }{
                \sin^2\left(\frac{\rotanglesymbol}{2}\right)
                \cdot \cos^2\left(\left(\frac{1}{2}+\repnum\right)\rotanglesymbol\right)
            } & \elsetext
        \end{cases} \, .
    \end{align*}
\end{lemma}
\begin{proof}
    Notice, that
    \begin{align*}
        \contractionof{\probofat{0}{\shortcatvariables},\bencodingofat{\formulaof{\selindex}}{\headvariableof{\selindex},\shortcatvariables}}{\headvariableof{\selindex}}
        = \coloredmatrixof{
            \left(\cosof{\frac{\rotanglesymbol}{2}}\right)^2 \\
            \left(\sinof{\frac{\rotanglesymbol}{2}}\right)^2
        }{\headvariableof{\selindex}}
        = \coloredmatrixof{
            1-\meanparamofat{0}{\indexedselvariable} \\
            \meanparamofat{0}{\indexedselvariable}
        }{\headvariableof{\selindex}}
    \end{align*}
    and
    \begin{align*}
        \contractionof{\probofat{\repnum}{\shortcatvariables},\bencodingofat{\formulaof{\selindex}}{\headvariableof{\selindex},\shortcatvariables}}{\headvariableof{\selindex}}
        =
        \coloredmatrixof{
            \left(\cosof{\left(\frac{1}{2}+\repnum\right)\rotanglesymbol}\right)^2 \\
            \left(\sinof{\left(\frac{1}{2}+\repnum\right)\rotanglesymbol}\right)^2
        }{\headvariableof{\selindex}}
        \coloredmatrixof{
            1-\meanparamofat{\repnum}{\indexedselvariable} \\
            \meanparamofat{\repnum}{\indexedselvariable}
        }{\headvariableof{\selindex}}
    \end{align*}
    If $\meanparamofat{\repnum}{\indexedselvariable}=1$ (that is $\sinof{\left(\frac{1}{2}+\repnum\right)\rotanglesymbol}=1$), the amplitude amplification amounts to adding $\formulaof{\selindex}$ as a hard constraint.
    This is equal to $\hardlegset^{\repnum}=\hardlegset\cup\{\selindex\}$ and $\headindexof{\hardlegset}^{\repnum}=(\headindexof{\hardlegset},1)$.

    If $\meanparamofat{\repnum}{\indexedselvariable}\neq 1$ we choose $\canparamofat{\Delta}{\indexedselvariable}\geq 0$ as in the claim and have after some algebra
    \begin{align*}
        \frac{
            \left(\cosof{\frac{\rotanglesymbol}{2}}\right)^2
        }{
            \left(\cosof{\frac{\rotanglesymbol}{2}}\right)^2 + \canparamofat{\Delta}{\indexedselvariable} \cdot \left(\sinof{\frac{\rotanglesymbol}{2}}\right)^2
        } = \left(\cosof{\left(\frac{1}{2}+\repnum\right)\rotanglesymbol}\right)^2
    \end{align*}
    and therefore
    \begin{align*}
        \contractionof{\probofat{\repnum}{\shortcatvariables},\bencodingofat{\formulaof{\selindex}}{\headvariableof{\selindex},\shortcatvariables}}{\headvariableof{\selindex}}
        = \normalizationof{
            \expof{\canparamofat{\Delta}{\indexedselvariable}\cdot\formulaofat{\selindex}{\canparam}}
            \probofat{0}{\shortcatvariables},\bencodingofat{\formulaof{\selindex}}{\headvariableof{\selindex},\shortcatvariables}
        }{\headvariableof{\selindex}} \, .
    \end{align*}
    Thus, the rotation of the q-sample corresponds with an increase of the canonical parameter $\canparamofat{0}{\indexedselvariable}$ by $\canparamofat{\Delta}{\indexedselvariable}$.
\end{proof}

\subsubsection{Maximum Entropy}

% Q-sample
When $\left(\frac{1}{2}+\repnum\right)\rotanglesymbol\leq \frac{\pi}{2}$ then all amplitudes remain positive, and $\qstateofat{\repnum}{\shortcatvariables}$ is a q-sample.
What is more, we can show that the property of maximum entropy with respect to a statistic containing the amplified formula.

\begin{theorem}
    \label{the:rotationKeepsMaxEntropy}
    Let $\qstateof{0}$ be the q-sample of a maximum entropy distribution with respect to the uniform base measure and the statistic $\formulaset$ containing the amplified formula, and let $\repnum\in\nn$ be a rotation number such that $\left(\frac{1}{2}+\repnum\right)\rotanglesymbol\leq\frac{\pi}{2}$.
    Then also $\qstateof{\repnum}$ is the q-sample of a maximum entropy distribution.
\end{theorem}
%We show \theref{the:rotationKeepsMaxEntropy} later based on a more technical description of the rotated state, which we show as the next lemma.
\begin{proof}%[Proof of \theref{the:rotationKeepsMaxEntropy}]
    By \lemref{lem:rotationCanParamIncrease} the amplitude amplified state is an elementary \ComputationActivationNetwork{} and thus a maximum entropy distribution.
        %   Since the measurement distribution of any quantum state in the plane spanned by $\qstateof{\goodstatesymbol}$ and $\qstateof{\badstatesymbol}$ is an elementary \ComputationActivationNetwork{}.
        %   When the assumption on $\repnum$ is met, we further have positive real amplitudes and therefore a q-sample.
        %  Since each amplitude amplification can be understood as an change of mean parameters.
        %  To be more precise, we characterize the change of canonical parameters in the lemma below.
\end{proof}

%
Notice, that we can only increase the canonical parameter by amplitude amplification.
Decreasing could be done by amplitude amplification on $\lnot\formulaof{\selindex}$ instead.



\subsection{Moment Matching} %based preparation of \CompActNets{}

We now investigate which \ComputationActivationNetworks{} can be prepared directly with amplitude amplification.
While we already have studied sparse representations of the computation network parts, we here apply amplitude amplification by Grover rotations as an activation mechanism.

We orient on the iterative proportional fitting algorithm, which is used to estimate the canonical parameters in exponential families  \cite{wainwright_graphical_2008} and more generally of \ComputationActivationNetworks{} \cite{goessmann_tensor-network_2025}.

The algorithm loops over the moment estimation and moment amplification subroutines, which have been introduced in the sections above.
The idea is to first compute the current moments of a distribution and second amplify one.



%This change of mean parameters can be understood by increasing the canonical parameter to $\exformula$ by
%\begin{align*}
%    \canparamof{\repnum} - \canparamof{0}
%    = \frac{
%        \cos^2\left(\frac{\rotanglesymbol}{2}\right) \left(1-
%                                                         \cos^2\left(\left(\frac{1}{2}+\repnum\right) \rotanglesymbol \right)
%        \right)
%    }{
%        \sin^2\left(\frac{\rotanglesymbol}{2}\right)
%        \cdot \cos^2\left(\left(\frac{1}{2}+\repnum\right)\rotanglesymbol\right)
%    }
%\end{align*}
%we here allow for $\canparam=\infty$, which corresponds to the case of a hard activation to $\exformula$.

We understand this rotation as a moment matching operation to the formula $\exformula$, see \algoref{alg:QCMM}.

\begin{algorithm}
    \caption{Quantum Circuit Moment Matcher}\label{alg:QCMM}
    \begin{algorithmic}
        \Require Formulas $\formulaofat{\seldim}{\shortcatvariables}$ for $\selindexin$, vector $\datameanat{\selvariable}$ of mean parameters
        \Ensure Circuit $\unitarysymbol$ preparing a state (when applied on the ground state) which measurement distribution matched $\datameanat{\selvariable}$\\
        \hrule
        \State $U\algdefsymbol (\paulixsymbol\circ\hgate)\left[\avariableof{\insymbol},\avariableof{\outsymbol}\right] \otimes
        \left(\bigotimes_{\catenumeratorin}\hgateat{\catvariableof{\catenumerator,\insymbol},\catvariableof{\catenumerator,\outsymbol}}\right)
        $
        \While{Convergence criterion not met}
%        \For{$\selindexin$}
            \State Estimate $\meanparamat{\selvariable}$,
            \begin{itemize}
                \item either by particle bases inference
                \item or by quantum counting (measuring the eigenvalue of the Grover operator), see \secref{sec:amplitudeEstimation}
            \end{itemize}
            \State Select an $\selindexin$, which moment has to be updated (e.g. by comparison of $\meanparamat{\selvariable}$ with $\datameanat{\selvariable}$)
            \State Choose whether to amplify $\formulaof{\selindex}$ or $\lnot\formulaof{\selindex}$ based on whether the estimated mean is smaller than the target mean
            \State Compute optimal $\repnum$ to match $\datameanat{\selvariable}$
            \State Extend Circuit $\unitarysymbol$ by $\repnum$ alternations of
            \begin{itemize}
                \item $\qcbencodingof{\formulaof{\selindex}}$ (effectively an reflection across the subspace spanned by one-hot encoded models)
                \item reflections across the current state (by $U(I-2e_0e_0)U$)
            \end{itemize}
        \EndWhile
    \end{algorithmic}
\end{algorithm}


% Outook
We can use \algoref{alg:QCMM} as a preparation algorithm of a q-sample before doing ancilla augmentation as in the main section.
Note that is important to keep track of the canonical parameters in the preparation.
When activating the statistic qubits, only the difference of the already prepared canonical parameters and the target canonical parameters has to be used (importance sampling interpretation).





\section{Extension: Sampling from proposal distributions}

We can prepare basis circuit encodings to selection augmented formulas, in this way introducing formula selecting networks.

\textbf{Idea for an inductive reasoning scheme:} Prepare a q-sample from the empirical distribution and the current distribution.
Then prepare the basis circuit encodings, where the selection variables are shared and the distributed variables assigned to the prepared samples.
Now, the ancilla qubits can be designed to $\onehotmapof{1}$ and $\onehotmapof{0}$ accordingly.
The rejection sampling scheme on both ancillas being $1$ and the measurement of $\selvariable$ prepares then the distribution
\begin{align*}
    \normalizationof{
        \contractionof{\empdistributionwith, \sencmlnstatwith}{\selvariable}, \contractionof{\currentdistributionwith}{\selvariable}
    }{\selvariable}
\end{align*}
That is, the probability of selecting $\selindex$ is proportional to
\begin{align*}
    \datameanat{\indexedselvariable} \cdot (1-\currentmeanat{\indexedselvariable})
\end{align*}
and thus prefers formulas, which have a large empirical mean, but a small current mean.

\textbf{Open Question:} Since the distribution is "similar" to $\expof{\datameanat{\indexedselvariable}-\currentmeanat{\indexedselvariable})}$ (terms appear in Taylor of first order), can we tune the distribution with an inverse temperature parameter $\beta$?

